\begin{Application}[][10 min]
Considérons les fonctions $f$ et $g$ définies sur $\mathbb{R}$ par :
 $f(x) = \cos(x) \quad \text{et} \quad g(x) = \sin(x) $
\begin{enumerate}
    \item Étudier la parité de $f$ et $g$.
    \item Calculer $f(x+2\pi)$ et $g(x+2\pi)$ ; déduire.
    \item Préciser puis compléter les courbes de $f$ et $g$ représentées ci-dessous.
\end{enumerate}

\begin{center}
    % --- PREMIER GRAPHIQUE : COSINUS ---
    \begin{tikzpicture}[x=0.9cm, y=1cm]
        % Axes
        \draw[->, >=stealth, thick] (-7.5,0) -- (7.5,0) node[below right] {$x$};
        \draw[->, >=stealth, thick] (0,-1.5) -- (0,1.5) node[left] {$y$};
        \node[below left] at (0,0) {$O$};

        % Lignes horizontales pointillées pour les limites (1 et -1)
        \draw[dashed, gray!80] (-7.5,1) -- (7,1);
        \draw[dashed, gray!80] (-7.5,-1) -- (7,-1);
        \node[left] at (0,1) {$1$};
        \node[left] at (0,-1) {$-1$};

        % Graduations de l'axe des abscisses en fonction de pi
        \foreach \k/\label in {
            -4/-2\pi, -3/-\frac{3\pi}{2}, -2/-\pi, -1/-\frac{\pi}{2},
             1/\frac{\pi}{2}, 2/\pi, 3/\frac{3\pi}{2}, 4/2\pi
        } {
            \draw (\k*pi/2, 0.1) -- (\k*pi/2, -0.1) node[below, font=\footnotesize] {$\label$};
        }

        % Tracé de la courbe (ébauche Cosinus)
        \draw[very thick, PrimaryColor, domain=0:pi, samples=100] plot (\x, {cos(\x r)});
        
        % Titre (optionnel)
        %\node[above, font=\bfseries] at (0, 1.8) {Fonction Cosinus};
    \end{tikzpicture}
    \hfill
    % --- DEUXIÈME GRAPHIQUE : SINUS ---
    \begin{tikzpicture}[x=0.9cm, y=1cm]
        % Axes
        \draw[->, >=stealth, thick] (-7.5,0) -- (7.5,0) node[below right] {$x$};
        \draw[->, >=stealth, thick] (0,-1.5) -- (0,1.5) node[left] {$y$};
        \node[below left] at (0,0) {$O$};

        % Lignes horizontales pointillées pour les limites (1 et -1)
        \draw[dashed, gray!80] (-7.5,1) -- (7,1);
        \draw[dashed, gray!80] (-7.5,-1) -- (7,-1);
        \node[left] at (0,1) {$1$};
        \node[left] at (0,-1) {$-1$};

        % Graduations de l'axe des abscisses en fonction de pi
        \foreach \k/\label in {
            -4/-2\pi, -3/-\frac{3\pi}{2}, -2/-\pi, -1/-\frac{\pi}{2},
             1/\frac{\pi}{2}, 2/\pi, 3/\frac{3\pi}{2}, 4/2\pi
        } {
            \draw (\k*pi/2, 0.1) -- (\k*pi/2, -0.1) node[below, font=\footnotesize] {$\label$};
        }

        % Tracé de la courbe (ébauche Sinus)
        \draw[very thick, PrimaryColor, domain=0:pi, samples=100] plot (\x, {sin(\x r)});
        
        % Titre (optionnel)
        %\node[above, font=\bfseries] at (0, 1.8) {Fonction Sinus};
    \end{tikzpicture}
\end{center}

\end{Application}
